% Generated draft: requires mathematical and source review.
\documentclass[11pt]{article}
\usepackage[margin=25mm]{geometry}
\usepackage{fontspec}
\setmainfont{Latin Modern Roman}
\usepackage{amsmath,amssymb,mathtools,unicode-math}
\setmathfont{Latin Modern Math}
\usepackage{enumitem,needspace,xurl,hyperref}
\setlength{\parindent}{0pt}
\setlength{\parskip}{5pt}
\setlength{\emergencystretch}{3em}
\newcommand{\N}{\mathbb{N}}
\newcommand{\Z}{\mathbb{Z}}
\newcommand{\Q}{\mathbb{Q}}
\newcommand{\R}{\mathbb{R}}
\newcommand{\C}{\mathbb{C}}
\newcommand{\F}{\mathbb{F}}
\newcommand{\A}{\mathbb{A}}
\newcommand{\PP}{\mathbb{P}}
\newcommand{\E}{\mathbb{E}}
\newcommand{\eps}{\varepsilon}
\newcommand{\dd}{\,\mathrm{d}}
\begin{document}
\section*{UnsolvedMath ID: 20000403}
\textbf{Source problem code:} catalogue entry
\subsection{Definitions and mathematical statement}
Brian Conrey asks to carry out ``(the other) Levinson's method'' for function fields. The original does not identify which Levinson argument, $L$-function family, mollifier, conductor/genus limit or intended zero-proportion conclusion is meant. The parenthetical is a genuine historical ambiguity.

The attached unreviewed aid treats a numerator-only analogue. Let $q>1$ be a prime power and $F(z)$ a nonconstant self-inversive complex polynomial, so $F(z)=\omega z^N\overline{F(1/\overline z)}$ with $|\omega|=1$, $N\ge1$, and nonzero constant term. Put $A(s)=F(q^{1/2-s})$ and choose the half-open period strip
\[
S_q=\{s\in\C:0\le\operatorname{Im}s<2\pi/\log q\}.
\]
The aid attributes to Cohn's theorem the multiplicity-counted equality between zeros of $A$ with $\operatorname{Re}s>1/2$ and zeros of $A'$ with $\operatorname{Re}s<1/2$. Self-inversivity then gives
\[
\#\{s\in S_q:A(s)=0,\ \operatorname{Re}s\ne1/2\}
=2\#\{s\in S_q:A'(s)=0,\ \operatorname{Re}s<1/2\}.
\]
Critical-line zeros are omitted from both open-half-plane counts, and the half-open convention avoids periodic duplication.

For a smooth projective geometrically connected curve $C/\F_q$ of positive genus, write $Z_C(u)=P_C(u)/((1-u)(1-qu))$. Taking $F(z)=P_C(q^{-1/2}z)$ gives $A(s)=P_C(q^{-s})$. The aid calls the resulting equivalence a numerator Speiser criterion. It makes no claim for the derivative of the full rational zeta function, whose denominator changes the derivative. The genus-zero constant-numerator case must be excluded from a derivative-zero count, because its derivative is identically zero.

Separately the aid reports that a fixed-function, finite-mollifier height mean is just a finite Parseval norm. For a finite polynomial $B(s)=\sum_j b_jq^{-js}$ and real $\sigma$, its period mean is $\sum_j|b_j|^2q^{-2j\sigma}$; longer height intervals repeat the same period. It therefore says that a nontrivial asymptotic Levinson-proportion implementation must vary conductor, degree or genus. The source does not supply that growing-family implementation or prove a new general zero-proportion result.
\subsection{Short English statement}
Specify which Levinson method and growing function-field family are intended. The aid offers a finite-period numerator derivative criterion and a Parseval calculation, keeping both separate from the full-zeta derivative and the missing family asymptotics.
\subsection{Sources}
\begin{itemize}
\item Brian Conrey, Question 25, 2014 AIM \emph{Arithmetic statistics over finite fields and function fields}: \url{https://aimath.org/pastworkshops/arithstatffieldproblems.pdf}. Canonical location: \texttt{aim-algebraic-number-theory-notes.json}, notes[54]. The original is labeled exact; remarks and literature are empty. No PDF was consulted.
\item The attached aid names Cohn's theorem and a Speiser-type criterion but supplies no complete bibliography or exact cited theorem. Its proposed identity has low novelty confidence and is explicitly machine generated and not peer reviewed.
\item Local registry: \texttt{data/unsolvedmath/problems.json}, permanent catalogue entry, exact source code \texttt{catalogue entry}.
\end{itemize}
\end{document}
